\section{Introduction}
Procedural memory lets an agent reuse a successful method without deriving it again. Yet success is conditional. A record may have been complete before reaching a parser; a cache may have belonged to one request; a diagnostic object may have contained only public fields. If the target changes these assumptions, a memory can accurately describe the source task while providing an incomplete guide to the new one.

Family F08 makes the problem concrete. In the source context, a status value is a decimal integer, so interpolating it into HTML cannot introduce markup. The target task changes that value to caller controlled text. Direct reuse still renders a status message, but fails the focal injection witness; a correct adaptation escapes the text before interpolation. The memory remains accurate about its source, while its applicability changes at the target boundary.

\begin{figure}[t]
\centering\includegraphics[width=\linewidth]{figures/design_overview.pdf}
\caption{The controlled intervention. A procedure is valid under an explicit source assumption. The target changes that assumption, then varies only the supplied memory condition. Functionality and the family specific focal witness are observed separately.}
\label{fig:design}
\end{figure}

We ask: \emph{when a coding agent receives a source procedure in a target context where its assumptions no longer hold, how does that memory change functional completion and focal security failures?} A second question asks whether a generic applicability reminder changes this behavior. We supply a frozen correct memory, a matched irrelevant memory, correct memory with the reminder, or no memory (Figure~\ref{fig:design}). The target task and changed assumption remain fixed. The original hypothesis predicted that correct memory would increase functional completions that violate the target's focal property. The estimates instead have positive, zero, and negative signs across the six configurations, which is inconsistent with a simple universal harmful transfer account in this cohort.

The study contributes (1) a controlled benchmark of thirteen security relevant applicability shifts; (2) a six configuration experiment that separates functionality from focal witness failure; and (3) implementation and visible trace evidence that distinguishes unchecked reuse, target appropriate protection, and failure to complete. The shifts span caller controlled paths, deferred resource lifetimes, cross user caches, origin changes, repeated control fields, concurrency invariants, and approximate authorization evidence. Together, the results show heterogeneous effects of correct memory and a configuration specific response to a generic applicability reminder.

All targets contain a context shift, and memories are supplied directly. The intervention therefore estimates the effect of supplied memory in those targets; it does not identify the effect of context shift or evaluate a retrieval system. Comparisons are made within each model and agent configuration.
